\subsection{Haces de $p$-diferenciales y diferencial exterior.}
\label{IV.16.6}

\begin{env}[16.6.1]
\label{IV.16.6.1}
Sea $f:X\to S$ un morfismo de preesquemas. Se llama \emph{haz de
$p$-diferenciales de $X$ relativo a $S$} ($p$ entero) a la
\emph{$p$-ésima potencia exterior} (0, 4.1.5) del $\mathcal{O}_X$-módulo
$\Omega_{X/S}^1$, denotada por
\[
  \Omega_{X/S}^p=\bigwedge^p(\Omega_{X/S}^1).
  \tag{16.6.1.1}
  \label{IV.16.6.1.1}
\]
Se tiene, por tanto, $\Omega_{X/S}^0=\mathcal{O}_X$, y $\Omega_{X/S}^p=0$ para $p<0$;
los $\Omega_{X/S}^p$ son los componentes homogéneos del álgebra exterior
sobre $\Omega_{X/S}^1$
\[
  \Omega_{X/S}^{\bullet}=\bigwedge(\Omega_{X/S}^1)
  =\bigoplus_{p\in\mathbf{Z}}\bigwedge^p(\Omega_{X/S}^1),
  \tag{16.6.1.2}
  \label{IV.16.6.1.2}
\]
que es, por tanto, una $\mathcal{O}_X$-álgebra graduada cuasicoherente,
anticonmutativa y cuyos elementos de grado 1 son de cuadrado nulo. Para todo
abierto afín $U$ de $X$, se tiene
\[
  \Gamma(U,\Omega_{X/S}^{\bullet})
  =\bigwedge(\Gamma(U,\Omega_{X/S}^1)),
\]
donde $\Gamma(U,\Omega_{X/S}^1)$ se considera como
$\Gamma(U,\mathcal{O}_X)$-módulo.

Cuando $S=\operatorname{Spec}(A)$ y $X=\operatorname{Spec}(B)$ son
afines, siendo $B$ una $A$-álgebra, se tiene (0, 4.1.5)
$\Omega_{X/S}^p=(\Omega_{B/A}^p)^{\sim}$, poniendo
$\Omega_{B/A}^p=\bigwedge^p\Omega_{B/A}^1$.
\end{env}

\begin{theorem}[16.6.2]
\label{IV.16.6.2}
Existe un único endomorfismo $d$ del haz de grupos aditivos
$\Omega_{X/S}^{\bullet}$ que tiene las propiedades siguientes:
\begin{enumerate}
\item[(i)] $d\circ d=0$.
\item[(ii)] Para todo abierto $U$ de $X$ y toda sección
$f\in\Gamma(U,\mathcal{O}_X)$, se tiene $df=d_{X/S}f$.
\oldpage[IV-4]{35}
\item[(iii)] Para todo abierto $U$ de $X$, todo par de enteros $p$, $q$ y
todo par de secciones
$\omega'_p\in\Gamma(U,\Omega_{X/S}^p)$,
$\omega''_q\in\Gamma(U,\Omega_{X/S}^q)$, se tiene
\[
  d(\omega'_p\wedge\omega''_q)
  =(d\omega'_p)\wedge\omega''_q
  +(-1)^p\omega'_p\wedge d\omega''_q.
  \tag{16.6.2.1}
  \label{IV.16.6.2.1}
\]
\end{enumerate}
Además, $d$ es un endomorfismo de $\psi^*(\mathcal{O}_X)$-módulos
graduados, de grado $+1$.
\end{theorem}

Supongamos probada la existencia del endomorfismo $d$. Para todo abierto afín
$U$ de $X$, toda sección de $\Omega_{X/S}^p$ sobre $U$ es (en virtud
de (i)) combinación lineal de un número finito de elementos de la forma
$g(df_1\wedge df_2\wedge\cdots\wedge df_p)$, donde $g$ y los $f_i$ son
secciones de $\mathcal{O}_X$ sobre $U$ (0, 20.4.7). Las condiciones (i)
y (iii) muestran entonces, por inducción sobre $p$, que necesariamente se tiene
\[
  d(g(df_1\wedge df_2\wedge\cdots\wedge df_p))
  =dg\wedge df_1\wedge df_2\wedge\cdots\wedge df_p.
  \tag{16.6.2.2}
  \label{IV.16.6.2.2}
\]
Esto prueba, por tanto, la \emph{unicidad} de $d$ y la última afirmación del teorema.
En virtud de esta propiedad de unicidad, para demostrar la existencia de $d$, se
puede reducir al caso en que $S=\operatorname{Spec}(A)$ y
$X=\operatorname{Spec}(B)$ son afines. Ahora bien (Bourbaki, \emph{Alg.}, cap.
III, 3\textsuperscript{a} ed., \S~10), para definir una $A$-antiderivación
$D$ de grado $+1$ de un álgebra exterior $\bigwedge(M)$ (donde $M$ es un
$B$-módulo y $B$ una $A$-álgebra), con valores en una
$A$-álgebra \emph{graduada anticonmutativa}
$C=\bigoplus_{n=0}^{\infty}C_n$, cuyos elementos de grado 1 son de cuadrado
nulo, basta elegir arbitrariamente una $A$-derivación $D_0$ de $B$
en $C_1$ y un $A$-homomorfismo $D_1$ de $M$ en $C_2$; existe entonces
una única $A$-antiderivación $D$ de $\bigwedge(M)$ en $C$
que coincide con $D_0$ en $B$ y con $D_1$ en $M$.

En el caso actual, $D_0$ es necesariamente igual a $d_{B/A}$ en virtud de
(ii); todo se reduce a ver, teniendo en cuenta (16.6.2.2), que existe un
$A$-homomorfismo $u$ de $\Omega_{B/A}^1$ en $\Omega_{B/A}^2$ tal que
se tenga
\[
  u(g.df)=dg\wedge df
  \tag{16.6.2.3}
  \label{IV.16.6.2.3}
\]
para cualesquiera $f$, $g$ en $A$; para ello bastará mostrar que
existe un $A$-homomorfismo
$v:B\otimes_A\Omega_{B/A}^1\to\Omega_{B/A}^2$ tal que
\[
  v(g.\omega)=dg\wedge\omega
  \tag{16.6.2.4}
  \label{IV.16.6.2.4}
\]
para $g\in B$ y $\omega\in\Omega_{B/A}^1$. Finalmente, como
$\Omega_{B/A}^1=\mathfrak{I}/\mathfrak{I}^2$ (donde
$\mathfrak{I}=\mathfrak{I}_{B/A}$ es el núcleo del homomorfismo canónico
$B\otimes_A B\to B$) y $\Omega_{B/A}^1$ está engendrado por los elementos
de la forma $g.df$, basta definir un $A$-homomorfismo
$w:B\otimes_A(B\otimes_A B)\to\Omega_{B/A}^2$ tal que
\[
  w(g'\otimes g\otimes f)=dg'\wedge(g.df)
  \tag{16.6.2.5}
  \label{IV.16.6.2.5}
\]
y tal que $w$ sea nulo sobre la imagen de $B\otimes_A\mathfrak{I}^2$. Ahora bien,
como el segundo miembro de (16.6.2.5) es $A$-trilineal en $g'$, $g$, $f$,
la existencia de $w$ que satisface (16.6.2.5) es inmediata. Como, por otra parte,
$\mathfrak{I}$ está engendrado por los elementos $1\otimes x-x\otimes1$
($x\in B$), se reduce a comprobar que, cuando
$z=(1\otimes x-x\otimes1)(1\otimes y-y\otimes1)$, se tiene
$w(g'\otimes z)=0$. Ahora bien, como
$z=1\otimes(xy)+(xy)\otimes1-x\otimes y-y\otimes x$, la fórmula (16.6.2.4)
muestra que
\oldpage[IV-4]{36}
basta ver que se tiene $d(xy)-x.dy-y.dx=0$, lo que expresa que $d$ es una
derivación.

Queda por probar que $d$ satisface la condición (i). Ahora bien, el cuadrado de una
antiderivación es una derivación (Bourbaki, \emph{loc. cit.}), y como
$\Omega_{B/A}^{\bullet}$ está engendrado por $\Omega_{B/A}^1$ como
$B$-álgebra, basta comprobar que $d(dz)=0$ para $z\in B$ y
$z\in\Omega_{B/A}^1$; en el primer caso, esto resulta de la fórmula
(16.6.2.3) con $g=1$; en el segundo, se puede reducir al caso en que $z=g.df$
con $f$, $g$ en $B$, y entonces, en virtud de (16.6.2.1) y (16.6.2.3), se tiene
\[
  d(d(g.df))=d(dg\wedge df)
  =(d(dg))\wedge(df)-(dg)\wedge(d(df))=0.
\]
\hfill Q.E.D.

\begin{definition}[16.6.3]
\label{IV.16.6.3}
La antiderivación $d$ definida en (16.6.2) (denotada también por $d_{X/S}$) se
llama \emph{diferencial exterior sobre $X$} (relativa a $S$).
\end{definition}

\begin{proposition}[16.6.4]
\label{IV.16.6.4}
Para todo cambio de base $g:S'\to S$, si se pone
$X'=X\times_S S'$, el homomorfismo canónico
\[
  \Omega_{X/S}^{\bullet}\otimes_S S'
  \longrightarrow\Omega_{X'/S'}^{\bullet}
  \tag{16.6.4.1}
  \label{IV.16.6.4.1}
\]
deducido del isomorfismo (16.5.9.1) es biyectivo. Además, si $s$ es una
sección de $\Omega_{X/S}^{\bullet}$ sobre un abierto $U$ de $X$,
$s\otimes1$ su imagen inversa, sección de $\Omega_{X'/S'}^{\bullet}$
sobre la imagen inversa $U'$ de $U$ en $X'$, se tiene
$d_{X'/S'}(s\otimes1)=d_{X/S}(s)\otimes1$.
\end{proposition}

La primera afirmación es inmediata, pues la formación del álgebra exterior
de un módulo conmuta con toda extensión del anillo de escalares. Para
probar la segunda, se puede, en virtud de (16.6.2.2), reducir al caso en que
$s\in\Gamma(U,\mathcal{O}_X)$, y en este caso la afirmación ya fue probada
(16.4.3.7).

\begin{env}[16.6.5]
\label{IV.16.6.5}
Supongamos que $\Omega_{X/S}^1$ es un $\mathcal{O}_X$-módulo localmente libre
de rango $n$ en un punto $x$, de modo que hay $n$ secciones
$s_i\in\Gamma(U,\mathcal{O}_X)$ tales que los $ds_i$ forman una base del
$\Gamma(U,\mathcal{O}_X)$-módulo $\Gamma(U,\Omega_{X/S}^1)$ (16.5.8).
Entonces, para todo entero $p\geq1$, los $p$-diferenciales
$ds_{i_1}\wedge ds_{i_2}\wedge\cdots\wedge ds_{i_p}$ (para
$i_1<i_2<\cdots<i_p$, elementos de $[1,n]$) forman una base de
$\binom{n}{p}$ elementos de $\Gamma(U,\Omega_{X/S}^p)$ sobre
$\Gamma(U,\mathcal{O}_X)$. Además, la fórmula (16.6.2.2) muestra que, para
toda sección $g\in\Gamma(U,\mathcal{O}_X)$, se tiene
\[
  d(g.ds_{i_1}\wedge ds_{i_2}\wedge\cdots\wedge ds_{i_p})
  =\sum_k(-1)^r\frac{\partial g}{\partial s_k}
  ds_{i_1}\wedge\cdots\wedge ds_{i_r}\wedge ds_k\wedge
  ds_{i_{r+1}}\wedge\cdots\wedge ds_{i_p}
  \tag{16.6.5.1}
  \label{IV.16.6.5.1}
\]
donde, en el segundo miembro, $k$ recorre el conjunto de los $n-p$ índices distintos de los
$i_h$, siendo $i_r$ el mayor índice $<k$.

Obsérvese que la relación $d(dg)=0$ para toda sección
$g\in\Gamma(U,\mathcal{O}_X)$ se expresa en la forma
\[
  D_i(D_jg)=D_j(D_ig)\qquad\text{para $i\neq j$};
\]
es decir, las derivaciones $D_i$ definidas en (16.5.7) conmutan dos a dos.
\end{env}
